\documentclass[12pt]{article}
\usepackage[utf8]{inputenc}
\usepackage[T1]{fontenc}
\usepackage{hyperref}
\title{Zapotec, Yatee Language Resources for K-12 Curriculum Development}
\author{Compiled by Claude Sonnet 5.5 (claude-sonnet-5-5) for Andr\'as Kornai}
\date{}
\begin{document}
\maketitle

\section{Introduction}
Yatee Zapotec (ISO 639-3: zty; Glottocode yate1242) is an Otomanguean language of Oaxaca, Mexico. It is considered stable, and Wikipedia reports about 5,000 speakers (2004). It is not known to be taught in schools. The searches found only a few resources.

\section{Primary Educational Resources}
No primers, textbooks or school programmes for Yatee Zapotec were found in search results.

\section{Linguistic and Reference Materials}
\begin{itemize}
\item San Francisco Yate\'e Zapotec Living Dictionary (Living Tongues Institute), searchable in English or Spanish: \url{https://livingdictionaries.app/san-francisco-yatee-zapo}
\item About page of the same Living Dictionary: \url{https://livingdictionaries.app/san-francisco-yatee-zapo/about}
\item Wikipedia article on Yatee Zapotec: \url{https://en.wikipedia.org/wiki/Yatee_Zapotec}
\item Glottolog entry for ISO code zty: \url{https://glottolog.org/resource/languoid/iso/zty}
\item Dictionary of the Zapotec Language, hosted by AILLA. It is a multi-variety dictionary and not specific to Yatee: \url{https://islandora-ailla.lib.utexas.edu/islandora/object/ailla%3A279600}
\end{itemize}

\section{Religious and Community Resources}
\begin{itemize}
\item Global Recordings Network page for Zapoteco, Yatee, listing audio and scripture resources: \url{https://globalrecordings.net/en/language/zty}
\end{itemize}

\section{Digital and Audio Resources}
The Living Dictionary and the Global Recordings Network page above are the digital and audio resources found. A New Testament recording from Faith Comes By Hearing was mentioned in the search summary, but no direct URL was seen, so it is not listed.

\section{Curriculum Development Context}
The language appears vitally used in its community, but the searches found no school materials. A curriculum would need to build on the Living Dictionary and on community-led documentation. Each village's variety should be checked, because Yatee and its neighbouring varieties differ.
\end{document}